(i) $X_{E} = (E - E_{0})^{2} + E(1 - E)$ is the system block of $(\widetilde E - \widetilde E_{0})^{2}$ for a purification (Proposition \claimref{synthesis.quasi-free-deposit-is-vacuum-plus-shot-noise}(ii)); the difference of two projectors has norm at most one, so $0 \le X_{E} \le 1$, and on the $g$-dimensional internal space at momentum $p$ this gives $x_{p} = \Tr X_{p} \le g$. Since $\omega_{p} \ge |p|$, the quanta with $\omega_{p} < k_{d}$ have $|p| < k_{d}$, and their density is at most $g\int_{|p| < k_{d}}d^{3}p/(2\pi)^{3} = gk_{d}^{3}/6\pi^{2}$. Write $\bar n = \int x_{p}$ for the density of quanta. On the support of $x$, $(\omega_{p} + k_{d})^{2} \le (\omega_{\max} + k_{d})(\omega_{p} + k_{d})$, so the integral is at most $(\omega_{\max} + k_{d})(\rho_{F} + k_{d}\bar n)$; and $k_{d}\bar n \le k_{d}\cdot gk_{d}^{3}/6\pi^{2} + \rho_{F}$, the first term from the quanta with $\omega_{p} < k_{d}$ and the second because $k_{d} x_{p} \le \omega_{p}x_{p}$ on the others. The factor $2$ is attained by quanta all at $\omega = k_{d}$. The check verifies $0 \le X_{E} \le 1$ and the trace formula on rational samples with non-commuting $E$ and $E_{0}$ (signatures by exact congruence), and the three inequalities on a cubic momentum grid for five occupation profiles, three masses and four values of $k_{d}$.

(ii) Insert (i) in the definition of $s$ with $\rho_{F} = f\rho$: $s \le \tfrac{3k_{d}^{3}}{2\pi^{2}\rho}(\omega_{\max} + k_{d})(2f + \eta)$. With $k_{d} = H_{d}/X$ and $M_{H} = 4\pi\rho/(3H_{d}^{3})$, $k_{d}^{3}/\rho = 4\pi/(3X^{3}M_{H})$, which gives the display; and $k_{d}^{4}/\rho = H_{d}^{4}/(X^{4}\rho) = 8\pi GH_{d}^{2}/(3X^{4})$ gives $\eta$. The check verifies both as rational identities with $\pi$ replaced by three rational values. Since $GH_{d}^{2} = (H_{d}/\bar M_{P})^{2}/8\pi$ is far below one in the window and $X > 10^{7}$, $\eta < 10^{-20}$ for any $g \le 10^{4}$.

(iii) The longitudinal part of the momentum density is that of the canonical one, since the symmetric momentum density differs from it by the curl of the spin density. Its cosine moment is the one-body observable with $f_{P} = \tfrac12(c\,p_{L} + p_{L}\,c)$, $c$ the multiplication by $\cos(\vec k\cdot\vec x)$ and $p_{L}$ the momentum along $\vec k$. As for the energy in the cited clause (iv), $f_{P}^{2} \le \tfrac12(c\,p_{L}^{2}\,c + p_{L}^{2})$, whose momentum-diagonal part is $\tfrac12\bigl(\tfrac14((p_{L} + k_{d})^{2} + (p_{L} - k_{d})^{2}) + p_{L}^{2}\bigr) = \tfrac34p_{L}^{2} + \tfrac14k_{d}^{2} \le \tfrac34(\omega_{p} + k_{d})^{2}$ because $|p_{L}| \le |p| \le \omega_{p}$. This is the bound the cited proof uses for the energy moment, so the cited clause (ii) gives the same shot term for the moments of the momentum density as for those of the energy density; dividing by $(\rho + p)^{2} = \tfrac{16}{9}\rho^{2}$ in place of $\rho^{2}$ gives the factor $\tfrac34$ in the amplitude.

(iv) The spectrum of a sum is bounded by the triangle inequality for standard deviations of the cosine and sine moments, as in the cited proof, so $\Delta_{\mathcal{R}}^{1/2} \le |\alpha|\Delta_{\delta}^{1/2} + |\beta|\Delta_{v}^{1/2}$, and the cited bound and (iii) give the display with $D = (|\alpha| + \tfrac34|\beta|)s^{1/2}$. For (FM), $\tfrac14 + 2X\cdot\tfrac34 = \tfrac14(1 + 6X)$. Solving $D^{2} \ge A$ with (ii): $s \ge 16A$ and $s \ge 16A/(1 + 6X)^{2}$ for (FR) and (FM), and $s \le \tfrac{2(2f + \eta)}{\pi X^{3}}(\omega_{\max} + k_{d})/M_{H}$ gives the first two thresholds. For (UE), $\tfrac32X^{2} + \tfrac34\cdot2X = \tfrac32X^{2}(1 + 1/X)$, and $s \ge 4A/(9X^{4}(1 + 1/X)^{2})$ gives the third. The coefficients $|\alpha|$ and $|\beta|$ of the three readings are those of Proposition \claimref{synthesis.descent-slice-constraints-and-the-three-readings-of-a-white-deposit} (its growing-mode amplitude in the radiation era, to leading order outside the Hubble length: a quarter of the contrast on a flat slice, $\tfrac32X^{2}$ times it on a slice of uniform expansion, and $2X$ times the velocity on every slice); they are not re-derived here, and which reading holds is a hypothesis.

(v) With $f \le 1$ the denominators are at most $2 + 10^{-20}$. The check evaluates the thresholds with $\pi > 333/106$ and $A = 2.0 \times 10^{-9}$ (the lower record value, Fact \claimref{synthesis.imported-planck-2018-lcdm-parameters}) at the stored lower bounds of $X = a_{d}H_{d}/a_{0}H_{0}$ (Theorem \claimref{synthesis.finite-range-descent-deposit-cannot-seed-cmb}(iii)): (FR) $3.0 \times 10^{15} < 8\pi AX^{3}/(2 + \eta) < 3.1 \times 10^{15}$ at $1$ MeV and above $9.6 \times 10^{60}$ at $10^{12}$ GeV; (FM) between $0.034$ and $0.035$ at $1$ MeV, above $5.0 \times 10^{13}$ at $10^{12}$ GeV, and increasing on a grid of $X$ between. $M_{H} = 1/(2GH_{d}) = 4\pi\bar M_{P}^{2}/H_{d}$ and $H_{d} = \pi\sqrt{g_{*}/90}\,T_{d}^{2}/\bar M_{P}$ (the cited theorem's form) give $(M_{H}/\bar M_{P})^{2} = 1440\,(\bar M_{P}/T_{d})^{4}/g_{*}$, more than $(6 \times 10^{43})^{2}$ at $1$ MeV and $(2 \times 10^{13})^{2}$ at $10^{12}$ GeV, and the products exceed $2.3 \times 10^{42}$ and $1.1 \times 10^{27}$. In reduced Planck masses the (FM) threshold is $\tfrac{8\pi A}{2f + \eta}\tfrac{X^{2}}{(1 + 6X)^{2}}\cdot XM_{H}$ with $XM_{H} = 4\pi\bar M_{P}^{2}a_{d}/k$, and $a_{d}$ decreases as $T_{d}$ rises (entropy conservation, as in the cited theorem) while $X^{2}/(1 + 6X)^{2}$ is within $10^{-7}$ of $\tfrac{1}{36}$ throughout, so the value at $10^{12}$ GeV is the least in the window. Under (FR) a white contrast with $\Delta_{\delta}^{2} = 16A$ at $X$ has $\Delta_{\delta}^{2} = 16AX^{3}$ at $k_{d} = H_{d}$. Under (UE), with $\pi < 22/7$, $A \le 2.2 \times 10^{-9}$, $f = 1$ and the stored $X > 3.2 \times 10^{20}$ at the pivot for $T_{d} = 10^{12}$ GeV, the threshold without the factor $(1 + 1/X)^{-2}$ is below $2.5 \times 10^{-30}$, and that factor is below one; and $(T_{d}/M_{H})^{2} = (T_{d}/\bar M_{P})^{2}(\bar M_{P}/M_{H})^{2} > 10^{-40}$ there, so $T_{d}/M_{H}$ exceeds the threshold by more than $10^{9}$.

The consequence combines (iv) and (v): they bound $\omega_{\max}$, the largest energy on the support of the occupation profile, so they show that some quantum lies above the threshold. The closing under (FR) and (FM) is conditional on Hypothesis (QP), which is stated in the result and is not among the stored premises; its motivation, that a single quantum with $G\omega^{2} \gg 1$ has a gravitational radius far above its wavelength and is not a quantum of a test field on the collar, is argued.
